\subsection{Sprint 1}

Tendo o escopo inicial pronto, as primeiras tasks foram definidas e divididas
(geralmente em duplas para facilitar com os AGES I) entre o time. Nesta sprint fiquei
inicialmente encarregado com uma task do backend, especificamente a implementação do login na parte de
autenticação do app.

Boa parte deste tempo acabei estudando e tentando entender o que devia ser
feito na minha task, devido a minha falta de conhecimento com as tecnologias
envolvidas. Como acabei pedindo ajuda muito tarde, dois colegas (Gabriel Lima e
Leonardo Berlatto, AGES II e III respectivamente) acabaram me auxiliando e me
mostrando o que deveria ser feito.
    
A primeira entrega para as stakeholders acabou sendo atrasada, e
consequentemente a task que eu estava envolvido. O AGES III que estava me
ajudando acabou terminando a atividade sozinho.
\\

\fbox{%
    \parbox{0.95\linewidth}{%
        \setlength{\parindent}{15pt}
            \itshape A maior lição que aprendi nesta sprint foi de pedir ajuda o mais cedo possível,
    acabei me sentindo culpado por atrasar o time com minha teimosia e timidez. Por
    mais que sou crítico da disciplina de AGES I, meus colegas do time sempre me
    ajudaram e eram extremamente prestativos. Eu não só poderia como deveria ter me
    expressado muito antes, e talvez assim, conseguiria ter feito mais pelo projeto.
    }%
}